
{\actuality} Актуальность темы диссертационного исследования обусловлена необходимостью
повышения эффективности процессов планирования и управления производством
в условиях многономенклатурного и мелкосерийного выпуска продукции.
Современное целлюлозно-бумажное производство представляет собой многостадийную
производственную систему, объединяющую выпуск бумажного полотна,
продольную резку, временное хранение полуфабрикатов и отгрузку продукции.
Планирование и управление в такой системе включают формирование производственных
планов, распределение заказов между машинами, расчёт раскроя,
составление календарных расписаний и контроль складских запасов.
В задачах раскроя и последующего формирования производственного расписания
для бумагоделательных машин и продольно-резательных станков
возникает необходимость одновременного учёта большого числа
технологических, временных и логистических ограничений, а также
дополнительных характеристик заказов, таких как сроки исполнения,
марка продукции, плотность, требования к кромке и числу слоёв.

Существующие математические модели задач раскроя и расписаний, как правило,
ориентированы либо на оптимизацию использования материала, либо на
минимизацию времени выполнения работ, при этом взаимосвязь между
процессами раскроя и формирования расписаний часто не учитывается.
В условиях реального производства это приводит к снижению качества
получаемых планов и увеличению потерь сырья.
Кроме того, детерминированные плановые решения не учитывают
случайные возмущения, возникающие при исполнении расписания:
отклонения фактических длительностей операций от плановых,
отказы и последующие ремонты оборудования.
Это обусловливает необходимость оценки устойчивости получаемых
планов методами имитационного моделирования.

В связи с этим разработка математических моделей и численных методов,
обеспечивающих согласованное решение задачи формирования
производственного расписания и раскроя с учётом расширенного набора ограничений,
а также их программная реализация и вычислительное исследование,
являются актуальной научной и прикладной задачей.

\ifsynopsis
Настоящая работа посвящена разработке математической модели
и численных методов решения задачи формирования расписаний и раскроя
в многомашинных производственных системах.
\else
Настоящая диссертационная работа направлена на развитие методов
математического моделирования и численных методов
задач раскроя и производственного планирования в сложных
технологических системах.
\fi

Классическая задача линейного раскроя и её обобщения исследованы в фундаментальных работах П.~Гилмора и Р.~Гомори \cite{gilmore1961linear, gilmore1963linear}, где предложены целочисленные формулировки и метод генерации столбцов.
Разработаны точные методы на основе декомпозиции Дантцига--Вульфа и ветвей и границ, а также широкий спектр приближённых методов, включая эвристики и метаэвристики (генетические алгоритмы, локальный поиск).
Для целлюлозно-бумажной промышленности предложены отраслевые модели, учитывающие потери при продольной резке рулонов, переналадки оборудования и ограничения по складам (А.~В.~Воронин \cite{воронин2000математические}, В.~А.~Кузнецов \cite{кузнецов2008многофункциональная}, Р.~В.~Воронов \cite{воронов2004диссертация}, А.~Р.~Урбан \cite{urban2015methods, урбан2015диссертация}).
Исследованы интегрированные постановки, объединяющие задачу раскроя с календарным планированием и планированием выпуска партий.
Однако существующие работы, как правило, либо ориентированы на оптимизацию использования материала, либо на минимизацию времени выполнения работ, при этом взаимосвязь между процессами раскроя и формирования расписаний часто учитывается в упрощённом виде.
Совместное рассмотрение нескольких ПРС и БРС, многослойных заказов, динамики складских запасов на уровне отдельных тамбуров, согласованного расписания и ограничения совместного размещения оборудования в единой модели в известных работах не встречается.

{\aim} данной работы является разработка математической модели,
численных методов и программных средств для решения задачи формирования
производственного расписания и раскроя на ЦБК
с учётом расширенного набора технологических и временных ограничений.

Для~достижения поставленной цели необходимо было решить следующие {\tasks}:
\begin{enumerate}[beginpenalty=10000] % https://tex.stackexchange.com/a/476052/104425
  \item провести анализ существующих подходов к моделированию и решению
  задач формирования производственных расписаний и раскроя;
  \item разработать формальную математическую постановку задачи,
  учитывающую взаимосвязь процессов раскроя и планирования работ;
  \item построить математическую модель задачи с учётом
  технологических, временных и логистических ограничений;
  \item разработать численный метод решения поставленной задачи,
  включая формирование модели производственного расписания;
  \item разработать имитационную модель исполнения производственного
  расписания, позволяющую оценивать устойчивость плановых решений
  к~случайным отклонениям длительностей операций и отказам оборудования;
  \item реализовать разработанные численные методы в виде программного комплекса;
  \item провести вычислительный эксперимент и проанализировать
  эффективность предложенных решений.
\end{enumerate}


{\novelty}
\begin{enumerate}[beginpenalty=10000] % https://tex.stackexchange.com/a/476052/104425
    \item Разработан метод математического моделирования совместного формирования и~оптимизации производственных планов и~расписаний в~целлюлозно-бумажной промышленности, в~которой решение представлено трёхкомпонентной структурой плана производства тамбуров, плана раскроя и~расписания.
    \item Предложена модель согласования стадий производства и~раскроя, обеспечивающая построение допустимого производственного расписания для цепочки БДМ, склад, ПРС/БРС с~учётом динамики складских состояний, доступности полуфабрикатов, сроков заказов и~технологических ограничений, включая требование размещения взаимодействующего оборудования (БДМ, ПРС/БРС и~складов) в~одном производственном помещении.
    \item Разработан численный метод решения оптимизационной задачи на~основе генетического алгоритма, в~котором поиск решения осуществляется в~пространстве допустимых интегрированных планов, а~модификация плана раскроя сопровождается согласованным пересчётом плана производства и~расписания с~учётом взаимосвязи стадий и~ресурсных ограничений.
    \item Разработана имитационная модель работы производственного расписания на~основе дискретно-событийного моделирования, учитывающая случайные отклонения фактических длительностей операций производства и~раскроя от~плановых, а~также отказы и~ремонты оборудования, и~позволяющая оценивать устойчивость полученного плана по~статистическим характеристикам времени завершения, числа просроченных заказов и~загрузки оборудования и~складов.
\end{enumerate}

{\influence} 
работы заключается в следующем. Теоретическая значимость исследования состоит в развитии методов математического моделирования и численных методов задач раскроя и производственного планирования в сложных технологических системах.
Предложенная трёхкомпонентная структура интегрированного плана, объединяющего план производства, план раскроя и расписание, расширяет существующие подходы к моделированию взаимосвязанных производственных процессов.
Модель расписания согласования стадий производства и раскроя с учётом динамики складских состояний и ограничений совместного размещения оборудования вносит вклад в развитие теории расписаний для многостадийных производственных систем.
Разработанная имитационная модель исполнения производственного расписания вносит вклад в развитие методов оценки устойчивости плановых решений в условиях стохастических возмущений.
Результаты работы могут быть использованы в дальнейших теоретических исследованиях задач интегрированного планирования в отраслях с многостадийной технологической цепочкой.

Практическая значимость работы подтверждается внедрением результатов
диссертационного исследования в практическую деятельность ООО
``Опти-Софт''.
Внедрены математическая модель задачи раскроя рулонов бумажного полотна
и формирования производственного расписания на ЦБК с~учётом технологических, временных и
логистических ограничений, а~также численные методы решения указанной
задачи, предназначенные для формирования согласованных планов
производства, нарезки и расписания работы оборудования,
а~также имитационная модель исполнения расписания, применяемая для
оценки устойчивости плановых решений к~случайным возмущениям.
Практическое применение результатов показало, что разработанная модель и
численные методы позволяют сократить время формирования производственных
планов и расписаний, обеспечить получение качественных плановых решений,
учитывать совокупность существенных производственных факторов, включая
параметры заказов, характеристики оборудования, ограничения складов и
требования к~последовательности выполнения операций, а~также снизить
трудоёмкость расчётов при сохранении требуемого качества планов
производства и раскроя.
Результаты работы могут быть использованы для повышения эффективности
процессов производственного планирования, в~том числе за~счёт сокращения
вычислительного времени при сохранении качества формируемых планов.

{\methods}
В работе использовались методы дискретной оптимизации,
теории расписаний, комбинаторной оптимизации,
математического и имитационного (дискретно-событийного) моделирования,
а также методы вычислительного эксперимента.

{\defpositions}
\begin{enumerate}[beginpenalty=10000] % https://tex.stackexchange.com/a/476052/104425
  \item Имитационная модель работы производственного расписания,
  предназначенная для оценки устойчивости плановых решений
  к~случайным возмущениям длительностей операций и отказам оборудования.
  \item Математическая модель задачи раскроя рулонов бумажного полотна 
  и формирования производственного расписания для 
  целлюлозно-бумажного комбината (ЦБК) с БДМ, каждая из которых имеет индивидуальные параметры выпуска,
  ПРС/БРС и складами, минимизирующая отходы материала, число раскроев и переналадок оборудования.
  \item Численный метод решения задачи раскроя рулонов 
  бумажного полотна и формирования производственного 
  расписания на основе генетического алгоритма.
  \item Комплекс программ для автоматизированного 
  планирования раскроя, формирования производственных 
  расписаний и имитационного моделирования их работы
  в целлюлозно-бумажной промышленности, 
  реализующий разработанные модели и численные методы.
\end{enumerate}
% В папке Documents можно ознакомиться с решением совета из Томского~ГУ
% (в~файле \verb+Def_positions.pdf+), где обоснованно даются рекомендации
% по~формулировкам защищаемых положений.

Вышеперечисленные результаты охватывают все три составляющие специальности 1.2.2~--- \emph{математическое моделирование}, \emph{численные методы} и \emph{комплексы программ}, а~также соответствуют направлениям исследований паспорта специальности:
п.~1 (новые методы математического моделирования)~--- трёхкомпонентная модель интегрированного плана производства, раскроя и~расписания и~модель согласования стадий производства и~раскроя;
п.~2 (разработка, обоснование и~тестирование вычислительных методов)~--- численные методы полного перебора и~генетического поиска;
п.~3 (реализация численных методов в~виде комплексов программ)~--- программный комплекс;
п.~6 (системы компьютерного и~имитационного моделирования)~--- имитационная модель исполнения расписания;
п.~8 (комплексные исследования с~применением математического моделирования и~вычислительного эксперимента)~--- интегрированная постановка задачи и~программная реализация;
п.~9 (постановка и~проведение численных экспериментов)~--- вычислительные эксперименты на~тестовых и~производственных экземплярах.

{\reliability} полученных результатов обеспечивается корректностью
математических построений, обоснованностью используемых методов,
а также результатами вычислительных экспериментов.
Полученные результаты согласуются с известными теоретическими положениями
и результатами, представленными в работах других авторов.

{\probation}
Основные результаты диссертационной работы докладывались и обсуждались
на научных конференциях и семинарах по тематике математического моделирования,
оптимизации и управления производственными системами.

{\contribution} Личный вклад автора заключается в формализации задачи,
разработке математической модели, численных методов решения,
программной реализации и проведении вычислительных экспериментов.
Все выносимые на защиту результаты получены автором самостоятельно.

\ifnumequal{\value{bibliosel}}{0}
{%%% Встроенная реализация с загрузкой файла через движок bibtex8. (При желании, внутри можно использовать обычные ссылки, наподобие `\cite{vakbib1,vakbib2}`).
    {\publications} Основные результаты по теме диссертации изложены
    в~XX~печатных изданиях,
    X из которых изданы в журналах, рекомендованных ВАК,
    X "--- в тезисах докладов.
}%
{%%% Реализация пакетом biblatex через движок biber
    \begin{refsection}[bl-author, bl-registered]
        % Это refsection=1.
        % Процитированные здесь работы:
        %  * подсчитываются, для автоматического составления фразы "Основные результаты ..."
        %  * попадают в авторскую библиографию, при usefootcite==0 и стиле `\insertbiblioauthor` или `\insertbiblioauthorgrouped`
        %  * нумеруются там в зависимости от порядка команд `\printbibliography` в этом разделе.
        %  * при использовании `\insertbiblioauthorgrouped`, порядок команд `\printbibliography` в нём должен быть тем же (см. biblio/biblatex.tex)
        %
        % Невидимый библиографический список для подсчёта количества публикаций:
        \phantom{\printbibliography[heading=nobibheading, section=1, env=countauthorvak,          keyword=biblioauthorvak]%
        \printbibliography[heading=nobibheading, section=1, env=countauthorwos,          keyword=biblioauthorwos]%
        \printbibliography[heading=nobibheading, section=1, env=countauthorscopus,       keyword=biblioauthorscopus]%
        \printbibliography[heading=nobibheading, section=1, env=countauthorconf,         keyword=biblioauthorconf]%
        \printbibliography[heading=nobibheading, section=1, env=countauthorother,        keyword=biblioauthorother]%
        \printbibliography[heading=nobibheading, section=1, env=countregistered,         keyword=biblioregistered]%
        \printbibliography[heading=nobibheading, section=1, env=countauthorpatent,       keyword=biblioauthorpatent]%
        \printbibliography[heading=nobibheading, section=1, env=countauthorprogram,      keyword=biblioauthorprogram]%
        \printbibliography[heading=nobibheading, section=1, env=countauthor,             keyword=biblioauthor]%
        \printbibliography[heading=nobibheading, section=1, env=countauthorvakscopuswos, filter=vakscopuswos]%
        \printbibliography[heading=nobibheading, section=1, env=countauthorscopuswos,    filter=scopuswos]}%
        %
        \nocite{*}%
        %
        {\publications} Основные результаты по теме диссертации изложены в~\arabic{citeauthor}~печатных изданиях,
        \arabic{citeauthorvak} из которых изданы в журналах, рекомендованных ВАК%
        \ifnum \value{citeauthorscopuswos}>0%
            , \arabic{citeauthorscopuswos} "--- в~периодических научных журналах, индексируемых Web of~Science и Scopus%
        \fi%
        \ifnum \value{citeauthorconf}>0%
            , \arabic{citeauthorconf} "--- в~тезисах докладов.
        \else%
            .
        \fi%
        \ifnum \value{citeregistered}=1%
            \ifnum \value{citeauthorpatent}=1%
                Зарегистрирован \arabic{citeauthorpatent} патент.
            \fi%
            \ifnum \value{citeauthorprogram}=1%
                Зарегистрирована \arabic{citeauthorprogram} программа для ЭВМ.
            \fi%
        \fi%
        \ifnum \value{citeregistered}>1%
            Зарегистрированы\ %
            \ifnum \value{citeauthorpatent}>0%
            \formbytotal{citeauthorpatent}{патент}{}{а}{}%
            \ifnum \value{citeauthorprogram}=0 . \else \ и~\fi%
            \fi%
            \ifnum \value{citeauthorprogram}>0%
            \formbytotal{citeauthorprogram}{программ}{а}{ы}{} для ЭВМ.
            \fi%
        \fi%
        % К публикациям, в которых излагаются основные научные результаты диссертации на соискание учёной
        % степени, в рецензируемых изданиях приравниваются патенты на изобретения, патенты (свидетельства) на
        % полезную модель, патенты на промышленный образец, патенты на селекционные достижения, свидетельства
        % на программу для электронных вычислительных машин, базу данных, топологию интегральных микросхем,
        % зарегистрированные в установленном порядке.(в ред. Постановления Правительства РФ от 21.04.2016 N 335)
    \end{refsection}%
    \begin{refsection}[bl-author, bl-registered]
        % Это refsection=2.
        % Процитированные здесь работы:
        %  * попадают в авторскую библиографию, при usefootcite==0 и стиле `\insertbiblioauthorimportant`.
        %  * ни на что не влияют в противном случае
        %\nocite{vakbib2}%vak        % TODO: заменить на реальные ключи из author.bib/registered.bib
        %\nocite{patbib1}%patent
        %\nocite{progbib1}%program
        %\nocite{bib1}%other
        %\nocite{confbib1}%conf
    \end{refsection}%
        %
        % Всё, что вне этих двух refsection, это refsection=0,
        %  * для диссертации - это нормальные ссылки, попадающие в обычную библиографию
        %  * для автореферата:
        %     * при usefootcite==0, ссылка корректно сработает только для источника из `external.bib`. Для своих работ --- напечатает "[0]" (и даже Warning не вылезет).
        %     * при usefootcite==1, ссылка сработает нормально. В авторской библиографии будут только процитированные в refsection=0 работы.
}

% При использовании пакета \verb!biblatex! будут подсчитаны все работы, добавленные
% в файл \verb!biblio/author.bib!. Для правильного подсчёта работ в~различных
% системах цитирования требуется использовать поля:
% \begin{itemize}
%         \item \texttt{authorvak} если публикация индексирована ВАК,
%         \item \texttt{authorscopus} если публикация индексирована Scopus,
%         \item \texttt{authorwos} если публикация индексирована Web of Science,
%         \item \texttt{authorconf} для докладов конференций,
%         \item \texttt{authorpatent} для патентов,
%         \item \texttt{authorprogram} для зарегистрированных программ для ЭВМ,
%         \item \texttt{authorother} для других публикаций.
% \end{itemize}
% Для подсчёта используются счётчики:
% \begin{itemize}
%         \item \texttt{citeauthorvak} для работ, индексируемых ВАК,
%         \item \texttt{citeauthorscopus} для работ, индексируемых Scopus,
%         \item \texttt{citeauthorwos} для работ, индексируемых Web of Science,
%         \item \texttt{citeauthorvakscopuswos} для работ, индексируемых одной из трёх баз,
%         \item \texttt{citeauthorscopuswos} для работ, индексируемых Scopus или Web of~Science,
%         \item \texttt{citeauthorconf} для докладов на конференциях,
%         \item \texttt{citeauthorother} для остальных работ,
%         \item \texttt{citeauthorpatent} для патентов,
%         \item \texttt{citeauthorprogram} для зарегистрированных программ для ЭВМ,
%         \item \texttt{citeauthor} для суммарного количества работ.
% \end{itemize}
% Счётчик \texttt{citeexternal} используется для подсчёта процитированных публикаций;
% \texttt{citeregistered} "--- для подсчёта суммарного количества патентов и программ для ЭВМ.

% Для добавления в список публикаций автора работ, которые не были процитированы в
% автореферате, требуется их~перечислить с использованием команды \verb!\nocite! в
% \verb!Synopsis/content.tex!.
